% =========================================================================
% Capítulo 1: Cinemática
% Tema 01: Coordenadas cilíndricas
% Problema 11: Disco en eje giratorio horizontal
% =========================================================================
\subsubsection{Disco en eje giratorio horizontal}

Un disco circular de radio $R$ se une a un eje giratorio de longitud $L$, como se muestra en la Figura \ref{fig:disco_eje_horizontal}. El eje gira respecto de la dirección horizontal con una velocidad angular constante $\Omega$ con respecto al suelo. El disco, a su vez, gira libremente en torno a su centro en el punto $O$. El ángulo $\theta$ describe la posición de un punto $P$ ubicado en el borde del disco relativo al centro del mismo. Utilizando los parámetros indicados, realice lo siguiente:

\begin{enumerate}
    \item Defina los marcos de referencia y los sistemas de coordenadas adecuados para describir el movimiento del sistema desde la perspectiva del suelo.
    \item Determine el vector posición del punto $P$ respecto al punto de apoyo $A$.
    \item Obtenga las expresiones para los vectores de velocidad angular y aceleración angular del eje y del disco.
    \item Calcule los vectores de velocidad y aceleración del punto $P$ respecto al suelo.
\end{enumerate}

\begin{figure}[!ht]
    \centering
    \begin{tikzpicture}[>=Latex, scale=1.8]
        % --- Parámetros Geométricos ---
        \def\L{3.2}          % Longitud del eje
        \def\R{1.1}          % Radio del disco
        \def\thetaAng{45}    % Ángulo theta para el punto P
        \def\BarThickness{0.06} % Grosor de la barra

        % --- Coordenadas ---
        \coordinate (A) at (0,0);
        \coordinate (O) at (\L,0);
        \coordinate (P) at ($ (O) + ({90-\thetaAng}:\R) $);

        % --- 1. Soporte en A ---
        \fill[pattern=north east lines] (-0.3, -0.4) rectangle (0.3, -0.3);
        \draw[thick] (-0.3, -0.3) -- (0.3, -0.3);
        \draw[thin] (A) -- (-0.15, -0.3);
        \draw[thin] (A) -- (0.15, -0.3);
        
        % Articulación A (Gris oscuro)
        \filldraw[fill=gray!80, draw=black] (A) circle (1.8pt) node[below left] {$A$};
    \shadedraw[inner color=white, outer color=gray!40, draw=black, thick] (O) circle (\R);
        % --- 2. Eje Giratorio (Barra L) ---
        \begin{scope}
            \shadedraw[inner color=white, outer color=gray!30, draw=black, rounded corners=2pt] 
                (0, -\BarThickness) rectangle (\L+0.1, \BarThickness);
        \end{scope}
        
        % Velocidad angular Omega
        \draw[->, thick] (1.2, 0.4) arc (60:300:0.2 and 0.4);
        \node[above] at (1.2, 0.4) {$\Omega$};
        
        % Cota de longitud L
        \draw[<->] (0, -0.6) -- (\L, -0.6) node[midway, fill=white] {$L$};
        \draw[dotted] (0, -0.1) -- (0, -0.7);
        \draw[dotted] (\L, -0.1) -- (\L, -0.7);

        % --- 3. Disco en O ---
        
        \draw[dashed] (O) -- ($(O) + (0, \R)$); % Referencia vertical
        
        % Radio R y Ángulo Theta
        \draw[->] (O) -- (P) node[midway, right] {$R$};
        \draw[->] ($(O) + (0, 0.5)$) arc (90:{90-\thetaAng}:0.5);
        \node at ($ (O) + (70:0.75) $) {$\theta$};

        % Articulación O (Gris oscuro)
        \filldraw[fill=gray!80, draw=black] (O) circle (1.2pt) node[above left, xshift=-2pt] {$O$};
        
        % Punto P
        \filldraw (P) circle (1.5pt) node[above right] {$P$};

    \end{tikzpicture}
    \caption{Representación esquemática de un disco de radio $R$ montado en un eje giratorio horizontal de longitud $L$.}
    \label{fig:disco_eje_horizontal}
\end{figure}

% --- SOLUCIÓN ---------------------------------------------------------------
\begin{solucion}{Disco en eje giratorio horizontal}
Lo primero que hacemos es definir los marcos de referencia, con los cuales trabajaremos durante el ejercicio.
\begin{itemize}
    \item $\mathcal{N}$ Suelo (Inercial) $\{\vec{E_x},\vec{E_y},\vec{E_z}\}$
    \item $\mathcal{A}$ Barra $AO$ $\{\vec{e_x},\vec{e_y},\vec{e_z}\}$
    \item $\mathcal{B}$ Disco  $\{\vec{e}_r,\vec{e}_\theta,\vec{e}_z\}$
\end{itemize}

\begin{center}
        \begin{tikzpicture}[>=Latex, scale=1.8]
        % --- Parámetros Geométricos ---
        \def\L{3.2}          % Longitud del eje
        \def\R{1.1}          % Radio del disco
        \def\thetaAng{45}    % Ángulo theta para el punto P
        \def\BarThickness{0.06} % Grosor de la barra

        % --- Coordenadas ---
        \coordinate (A) at (0,0);
        \coordinate (O) at (\L,0);
        \coordinate (P) at ($ (O) + ({90-\thetaAng}:\R) $);

        % --- 1. Soporte en A ---
        \fill[pattern=north east lines] (-0.3, -0.4) rectangle (0.3, -0.3);
        \draw[thick] (-0.3, -0.3) -- (0.3, -0.3);
        \draw[thin] (A) -- (-0.15, -0.3);
        \draw[thin] (A) -- (0.15, -0.3);
        
        % Articulación A (Gris oscuro)
        \filldraw[fill=gray!80, draw=black] (A) circle (1.8pt) node[below left] {$A$};
    \shadedraw[inner color=white, outer color=gray!40, draw=black, thick] (O) circle (\R);
        % --- 2. Eje Giratorio (Barra L) ---
        \begin{scope}
            \shadedraw[inner color=white, outer color=gray!30, draw=black, rounded corners=2pt] 
                (0, -\BarThickness) rectangle (\L+0.1, \BarThickness);
        \end{scope}
        
        % Velocidad angular Omega
        \draw[->, thick] (1.2, 0.4) arc (60:300:0.2 and 0.4);
        \node[above] at (1.2, 0.4) {$\Omega$};
        
        % Cota de longitud L
        \draw[<->] (0, -0.6) -- (\L, -0.6) node[midway, below] {$L$};
        \draw[dotted] (0, -0.1) -- (0, -0.7);
        \draw[dotted] (\L, -0.1) -- (\L, -0.7);

        % --- 3. Disco en O ---
        
        \draw[dashed] (O) -- ($(O) + (0, \R)$); % Referencia vertical
        
        % Radio R y Ángulo Theta
        \draw[->] (O) -- (P) node[midway, right] {$R$};
        \draw[->] ($(O) + (0, 0.5)$) arc (90:{90-\thetaAng}:0.5);
        \node at ($ (O) + (70:0.75) $) {$\theta$};

        % Articulación O (Gris oscuro)
        \filldraw[fill=gray!80, draw=black] (O) circle (1.2pt) node[above left, xshift=-2pt] {$O$};
        
        % Punto P
        \filldraw (P) circle (1.5pt) node[above right] {$P$};

        %aqui van los vectores
        \draw[->, thick, red] (0,0)--++(0.5,0) node[below] {$\vec{E}_y$};
        \draw[->, thick, red] (0,0)--++(0,0.5) node[above] {$\vec{E}_z$};
        
        %aqui van los vectores
        \draw[->, thick, blue] (\L,0)--++(0.7,0) node[below] {$\vec{e}_x$};
        \draw[->, thick, blue] (\L,0)--++(0,0.7) node[above] {$\vec{e}_y$};
        %aqui van los vectores
        \draw[->, thick, green] (\L,0)--++(45:0.7) node[above] {$\vec{e}_r$};
        \draw[->, thick, green] (\L,0)--++(-45:0.7) node[below] {$\vec{e}_\theta$};

    \end{tikzpicture}
\end{center}
\end{solucion}

